\documentclass{article}
\usepackage[utf8]{inputenc}
\usepackage[brazilian]{babel}
\usepackage[T1]{fontenc}    % Garante a formatação correta das fontes com acentos

\title{Revisão de Literatura: Arquitetura em Sistemas Embarcados para Experimentos Didáticos de Física}
\author{Victor Hugo Santos Bitencourt}
\date{\today}

\begin{document}

\maketitle

\section{Introdução}
Esse documento traz análise de referências de conteúdo e de definições de conceitos para serem utilizadas em monografia final.

\section{Referências que trazem relevância}
\begin{itemize}
\item{Formação técnico-profissional}
\begin{itemize}
\item Ao trabalhar com aparatos experimentais automatizados, que eles mesmos constroem ou modificam, alunos e professores são duplamente desafiados: por um lado, está em jogo a compreensão dos fenômenos investigados, por outro, a compreensão dos dispositivos, procedimentos e algoritmos, e do seu impacto sobre os resultados dos experimentos. [\dots] Com isso, aproxima-se o modo de se fazer física experimental, no ambiente didático, do modo como ela é realmente feita no âmbito profissional. \cite{rogerio}
\item Em tese, no contexto educacional, o cotejamento entre teoria e empiria seria uma
forma de proporcionar aos estudantes o ensejo de apreender a natureza dos fenômenos,
mas também desenvolver o senso crítico em relação aos procedimentos e aparatos usados
para construir conhecimento. Portanto, além dos fenômenos físicos, propriamente ditos,
está em jogo a questão do método, ou seja, dos caminhos que nos levam a quantificá-los
e explicá-los dessa ou daquela maneira. \cite{rogerio}
\end{itemize}
\item{Gestão de telemetria}
\begin{itemize}
\item Esse ambiente reduz o esforço com o monitoramento dos experimentos e o registro dos dados, liberando tempo para alunos e professores investigarem a natureza dos fenômenos e aprofundarem considerações sobre os métodos experimentais. \cite{rogerio}
\end{itemize}
\end{itemize}

\section{Referências que explicam conceitos}
\begin{itemize}
	\item Circuito RC
	\begin{itemize}
	\item No regime transitório, um circuito RC simples consiste em resistor, capacitor e fonte de tensão, onde o polo positivo do capacitor encontra-se conectado em série com o resistor, que por sua vez, é conectado ao polo positivo da fonte. Os polos negativos da fonte de tensão e do capacitor encontram-se conectados, criando uma diferença de potencial entre os terminais do capacitor. Inicialmente, a chave encontra-se aberta e o capacitor está descarregado, com tensão nula entre seus terminais. Ao se fechar a chave no instante t = 0, a tensão entre os terminais do resistor é a mesma da fonte e a tensão entre os terminais do capacitor é nula (Young; Freedman, 2009, p. 182). A partir desse instante, o processo de carga inicia-se no capacitor e a tensão entre seus terminais aumenta com o passar do tempo, na medida em que o capacitor acumula cargas positivas e cargas negativas nas suas placas condutoras, separadas por um isolante. A carga é mais rápida no início do processo e tornar-se mais lenta na proporção em que tensão no capacitor aproxima-se da tensão na fonte, até que essas tensões sejam praticamente iguais e o fluxo de corrente no circuito cesse. Logo, o capacitor estará totalmente carregado quando quando a tensão entre seus terminais for (praticamente) a mesma da fonte. A tensão sobre os terminais do capacitor, em função do tempo, é dada pela equação V = Vcc · (1 − e-t/RC), onde R é resistência, expressa em Ohms, C é a capacitância, expressa em Farads, e t é o tempo, em segundos, decorrido desde o início da carga. Logo (em condições ideais), a tensão da fonte é a assíntota da qual a tensão sobre o capacitor se aproxima, indefinidadmente (Young; Freedman, 2009, p.184) \cite{rogerio}
	\end{itemize}
\end{itemize}

\section{Referências que trazem estado da arte}
Fale sobre o tópico.

\section{Referências que embasam justificativas para a intervenção}
\begin{itemize}
	\item{Baixo Custo}
		\begin{itemize}
			\item As placas de desenvolvimento (ou prototipagem), assim como os transdutores (sensores e atuadores) podem hoje (2026) ser adquiridos com orçamento muito modesto. De fato, é possível automatizar muitos experimentos didáticos relevantes com investimento inferior a 200 reais. \cite{rogerio}
		\end{itemize}
	\item{Experimento do circuito RC}
		\begin{itemize}
			\item [\dots] envolvendo tópicos fundamentais da eletrodinâmica e da eletrostática, logaritmos e equações exponenciais, tanto no nível do ensino médio como no nível universitário. \cite{rogerio}
		\end{itemize}
	\item{Arduino}
		\begin{itemize}
			\item de longe, a escolha mais popular para automatizar o controle e aquisição de dados em experimentos didáticos, por causa do seu preço reduzido, ambiente de programação amigável e ampla disponibilidade de bibliotecas, referências de uso e tutoriais. O Arduino UNO R3 (UNO R3), usado no presente trabalho, atende bem aos requisitos de diversos experimentos didáticos de física (Martinazzo et al., 2014; Fetzner, 2015; Calderon, 2018).  \cite{rogerio}
			\item Um detalhe técnico importante, que contribui para a adequação do Uno R3 aos experimentos didáticos, é a linearidade e a confiabilidade na leitura e conversão de sinais analógicos, em comparação com outras placas de desenvolvimento (Monk, 2024). \cite{rogerio}
		\end{itemize}
	\item{Raspberry Pi}
		\begin{itemize}
		\item O baixo custo e a robustez fazem do Raspberry Pi uma excelente escolha para o ambiente de laboratório, onde o manuseio frequente por professores e estudantes expõe os equipamentos a maiores riscos de mal funcionamento. \cite{rogerio}
		\end{itemize}
\end{itemize}

\section{Referências que dão credibilidade às afirmações}
Fale sobre o tópico.

\section{Referências que trazem validação}
Fale sobre o tópico.

\section{Imagens de referências}
Adicione aqui imagens com referências.

\section{Conclusão}
Conclusão aqui.

\bibliographystyle{plain}
\bibliography{references}

\end{document}
